% META: {"id": "1NSI-TB-COURS-C4", "chapitre": "1NSI-TYPES-BASE", "type_objet": "cours", "capacites": ["P-BASE-04"], "mode_creation": "ex_nihilo", "sources_inspiration": [], "status": "verified"}
\section{Valeurs et opérateurs booléens}

\definition[D5]{
  Une \textbf{valeur booléenne} ne prend que deux valeurs : \lstinline{True} (vrai, souvent
  codé 1) ou \lstinline{False} (faux, souvent codé 0). Les trois opérateurs booléens de base
  sont \lstinline{and} (et), \lstinline{or} (ou) et \lstinline{not} (non).
}

\propriete[Table de vérité de \lstinline{and} et \lstinline{or}]{
  \begin{center}
  \begin{tabular}{cc|cc}
    $a$ & $b$ & $a\ \text{and}\ b$ & $a\ \text{or}\ b$ \\\hline
    Vrai & Vrai & Vrai & Vrai \\
    Vrai & Faux & Faux & Vrai \\
    Faux & Vrai & Faux & Vrai \\
    Faux & Faux & Faux & Faux \\
  \end{tabular}
  \end{center}
  \lstinline{and} est vrai seulement si \textbf{les deux} opérandes sont vrais ;
  \lstinline{or} est vrai dès qu'\textbf{au moins un} des deux opérandes est vrai.
}

% BEGIN-VERIFY
% table_and = {(a, b): a and b for a in (True, False) for b in (True, False)}
% assert table_and == {(True, True): True, (True, False): False, (False, True): False, (False, False): False}
% table_or = {(a, b): a or b for a in (True, False) for b in (True, False)}
% assert table_or == {(True, True): True, (True, False): True, (False, True): True, (False, False): False}
% END-VERIFY

\definition[D6]{
  Le \textbf{ou exclusif} (\lstinline{xor}, noté $\oplus$) est vrai lorsque \textbf{exactement
  un} des deux opérandes est vrai (mais pas les deux). En Python, on l'obtient avec
  \lstinline{a != b} pour des booléens, ou l'opérateur \lstinline{^} bit à bit sur des entiers.
}

\exemple{Le xor sert par exemple à additionner deux bits sans retenue : $0\oplus0=0$,
$0\oplus1=1$, $1\oplus0=1$, $1\oplus1=0$ (la retenue, elle, correspond à \lstinline{a and b}).}

\begin{codereference}{Addition d'un bit (demi-additionneur)}
\begin{python}
def demi_additionneur(a, b):
    """Additionne deux bits (0 ou 1) : renvoie (somme, retenue)."""
    somme = a ^ b
    retenue = a & b
    return somme, retenue

print(demi_additionneur(1, 1))  # (0, 1) : 1+1 = 10 en binaire
print(demi_additionneur(1, 0))  # (1, 0)
\end{python}
\end{codereference}

% BEGIN-VERIFY
% def demi_additionneur(a, b):
%     somme = a ^ b
%     retenue = a & b
%     return somme, retenue
% assert demi_additionneur(0, 0) == (0, 0)
% assert demi_additionneur(0, 1) == (1, 0)
% assert demi_additionneur(1, 0) == (1, 0)
% assert demi_additionneur(1, 1) == (0, 1)
% END-VERIFY

\propriete[Table de vérité d'une expression composée]{
  Pour dresser la table de vérité d'une expression booléenne à $n$ variables, on énumère les
  $2^n$ combinaisons possibles de valeurs des variables, et on calcule la valeur de
  l'expression pour chacune.
}

\exemple{Table de vérité de $(a\ \text{and}\ b)\ \text{or}\ (\text{not}\ a\ \text{and}\ c)$ :}

% BEGIN-VERIFY
% def expr(a, b, c):
%     return (a and b) or ((not a) and c)
% resultats = {(a, b, c): expr(a, b, c)
%              for a in (True, False) for b in (True, False) for c in (True, False)}
% assert resultats[(True, True, False)] is True
% assert resultats[(True, False, False)] is False
% assert resultats[(False, False, True)] is True
% assert resultats[(False, False, False)] is False
% assert len(resultats) == 8
% END-VERIFY

\erreurFrequente{Croire que \lstinline{and} et \lstinline{or} évaluent toujours leurs deux
opérandes. En Python, l'évaluation est \textbf{séquentielle} et \textbf{court-circuitée} :
dans \lstinline{a and b}, si \lstinline{a} vaut \lstinline{False}, \lstinline{b} n'est même
pas évalué (le résultat est déjà déterminé). Cela peut avoir des effets de bord si
\lstinline{b} contient un appel de fonction.}

% BEGIN-VERIFY
% compteur = {'appels': 0}
% def effet_de_bord():
%     compteur['appels'] += 1
%     return True
% resultat = False and effet_de_bord()
% assert resultat is False
% assert compteur['appels'] == 0  # effet_de_bord() n'a jamais ete appelee : court-circuit
% END-VERIFY

\alamachine{Écris une fonction \lstinline{table_verite(expr, n)} qui... en fait, dresse
« à la main » (avec des boucles \lstinline{for}) la table de vérité de l'expression
\lstinline{(a or b) and (not c)} pour les 8 combinaisons de \lstinline{a}, \lstinline{b},
\lstinline{c}. Vérifie ton résultat avec Python.}
